\documentclass[../DoAn.tex]{subfiles}
\begin{document}

% =========================
% CHƯƠNG 3 (ngắn gọn, tổng quan & định vị đóng góp)
% =========================

\section{Giới thiệu chương}
% TODO: viết thành văn xuôi (gạch đầu dòng định hướng)
\begin{itemize}
  \item Phát biểu bài toán tách fECG từ aECG và vai trò của tiền xử lý.
  \item Bố cục chương: nhóm phương pháp cổ điển (3.2), nhóm học sâu (3.3), nhận xét \& định vị đề tài (3.4).
\end{itemize}

\section{Nhóm phương pháp cổ điển}

\subsection{Bối cảnh và vai trò của tiền xử lý}
Trong đo ECG bụng (aECG), tín hiệu thu được luôn là hỗn hợp gồm mECG, fECG và nhiều nguồn gây nhiễu/hoạt động sinh học khác (như hoạt động cơ của mẹ, cử động thai, nhiễu môi trường). Do đó, \textit{tiền xử lý} là bước gần như bắt buộc trước khi áp dụng các thuật toán tách fECG. Tuy nhiên, tiền xử lý cũng có mặt trái: nếu lọc/khử nhiễu quá mạnh có thể làm mất các thành phần biên độ nhỏ nhưng có giá trị sinh lý (đặc biệt là sóng P và T), khiến tín hiệu sau tách bị “mượt” nhưng mất đặc trưng hình thái.

Trong các quy trình cổ điển, tiền xử lý thường nhằm loại (i) trôi đường nền (baseline wander), (ii) nhiễu lưới điện (power-line), và (iii) nhiễu tần cao. Một cấu hình điển hình là notch 50\,Hz để triệt nhiễu lưới, kết hợp lọc thông cao để xử lý trôi nền và lọc thông thấp để giảm nhiễu tần cao.

\subsection{Nhóm phương pháp lọc (Filtering-based)}
\subsubsection{Lọc thích nghi (Adaptive filtering)}
Lọc thích nghi là hướng tiếp cận cổ điển, tiêu biểu như LMS/NLMS và RLS \cite{fotiadou_enhancement_2018}. Trực giác của nhóm này là “học” bộ lọc để khử thành phần mECG/nhiễu dựa trên tiêu chuẩn tối ưu sai số. Điểm yếu quan trọng là thường cần \textit{tín hiệu tham chiếu} bổ sung để tách các thành phần trong aECG; đồng thời tính ổn định/phụ thuộc hội tụ khiến hiệu năng nhạy với lựa chọn tham số (ví dụ bước học của LMS).

\subsubsection{Kalman Filter (KF) và Extended Kalman Filter (EKF)}
KF thuộc họ lọc Bayes theo mô hình không gian trạng thái: dùng phép đo quan sát và mô hình động lực để ước lượng tuần tự trạng thái ẩn theo thời gian \cite{suganthy_enhancement_2019}. KF chuẩn giả định hệ tuyến tính; trong thực tế ECG và nhiễu thường phi tuyến/phi Gauss, nên EKF được dùng như một mở rộng: tuyến tính hóa (xấp xỉ) mô hình phi tuyến quanh kỳ vọng để vẫn áp dụng khung Kalman \cite{sarafan_novel_2022}. Ý nghĩa thực hành của KF/EKF trong bài toán fECG là: nếu mô hình hóa được động lực mECG hoặc ECG “sạch”, ta có thể ước lượng mECG rồi trừ khỏi aECG để lộ fECG (tương tự tinh thần của template subtraction nhưng “thích nghi theo thời gian”).

\subsubsection{Các kỹ thuật khử nhiễu thời--tần và phân rã tín hiệu}
Ngoài KF/EKF, nhiều kỹ thuật khử nhiễu cổ điển khác cũng thường xuất hiện trong hệ sinh thái ECG như làm trơn (Savitzky--Golay), khử nhiễu Wavelet, EMD/EEMD, Total Variation Denoising (TVD), và các biến thể dựa trên tính thưa (sparsity). Các nhóm này chủ yếu phục vụ mục tiêu tăng SNR trước khi bước vào tách nguồn hoặc phát hiện QRS.

\subsection{Nhóm phương pháp tách nguồn mù (Blind Source Separation -- BSS)}
\subsubsection{Independent Component Analysis (ICA)}
BSS giả định aECG là tổng tuyến tính của nhiều nguồn độc lập (fECG, mECG và nhiễu) \cite{cardoso_blind_1998}. ICA mô hình hóa hỗn hợp quan sát $X$ như $X = AS$ và tìm ma trận “giải trộn” $W$ để $Y = WX \approx S$. Thực tế triển khai ICA thường gồm các bước chuẩn hóa như \textit{centering} (đưa về trung bình 0), \textit{whitening} (làm trắng để các thành phần không tương quan và có phương sai chuẩn hóa), sau đó là bước lặp tối ưu dựa trên tiêu chuẩn phi-Gauss (ví dụ cực đại độ phi-Gauss, hợp lý cực đại, hoặc cực tiểu thông tin tương hỗ).

Các biến thể ICA hay dùng gồm JADE \cite{cardoso_blind_1993}, FastICA \cite{hyvarinen_independent_2000} và RobustICA \cite{zarzoso_robust_2010}. FastICA phổ biến vì đơn giản và hội tụ nhanh, phù hợp hướng “real-time” nhờ khả năng song song hóa và tìm từng thành phần độc lập. RobustICA sử dụng hàm tương phản kurtosis tổng quát và cơ chế bước học thích nghi; nhờ đó có thể giảm rủi ro kẹt tại nghiệm xấu và (theo mô tả cổ điển) bớt phụ thuộc một số tiền xử lý.

\subsection{Nhóm phương pháp khử mECG theo mẫu (Template Subtraction -- TS)}
\subsubsection{Template Subtraction cơ bản}
Ý tưởng TS là tái tạo mECG từ các nhịp của mẹ rồi trừ khỏi aECG. Quy trình điển hình: (i) phát hiện mQRS, (ii) cắt các chu kỳ mECG quanh đỉnh R của mẹ theo một cửa sổ thời gian xác định, (iii) lấy trung bình để tạo template mECG, (iv) lặp lại template theo thời gian để tái tạo mECG, và (v) trừ khỏi aECG để thu tín hiệu dư (residual) giàu thành phần fECG hơn. Nút thắt của TS là độ tin cậy của mQRS detection; khi fQRS chồng lấp thời gian với mQRS, việc phát hiện mQRS trở nên khó và kéo theo sai lệch lớn trong template, làm suy giảm hiệu quả tách fECG.

\subsubsection{Template Subtraction có hiệu chỉnh biên độ (TSc)}
Do hình thái mECG thay đổi theo thời gian, template trung bình có thể lệch so với nhịp thực tại từng thời điểm. Biến thể TSc khắc phục bằng cách nhân template với một hệ số tỉ lệ $\alpha$ để giảm sai khác giữa chu kỳ mECG thực và template; hệ số này có thể được chọn theo tiêu chuẩn tối thiểu hóa sai số bình phương (LMS) giữa nhịp và template đã co giãn \cite{martens_robust_2007}.

\subsection{Kết hợp TS và ICA (Hybrid TS--ICA)}
Sự kết hợp của các phương pháp có thể dẫn đến hiệu năng tốt hơn \cite{behar_non-invasive_2013}. Do mỗi phương pháp có điểm mạnh/điểm yếu khác nhau, các tổ hợp TS và ICA thường được dùng để tăng độ vững. Ba cấu hình tổ hợp thường gặp:
(1) TS--ICA: dùng TS để giảm thành phần mECG trội biên độ trước, rồi ICA tách phần còn lại;
(2) ICA--TS: ICA tách tương đối các thành phần, sau đó TS “dọn” phần mECG còn rò trong kênh fECG;
(3) ICA--TS--ICA: thêm một lượt ICA sau khi đã loại mECG nhằm làm sạch hơn các nguồn còn lại.

Một pipeline đánh giá điển hình cho tổ hợp TS/ICA thường đi qua: tiền xử lý (notch 50\,Hz, high-pass và low-pass), phát hiện mQRS bằng Pan--Tompkins để tạo template hoặc hỗ trợ chọn kênh mECG trong đầu ra ICA, sau đó phát hiện fQRS bằng Pan--Tompkins trên residual/đầu ra đã tách. Với ICA, việc chọn “kênh fECG” có thể dựa trên chỉ báo độ trơn nhịp (ví dụ đếm số lần biến thiên nhịp tức thời vượt ngưỡng trong các đoạn 1 phút), nhằm loại các kênh nhiễu có biến thiên nhịp bất thường.

\subsection{Nhận xét chung về giới hạn của các phương pháp cổ điển}
Nhìn tổng thể, các phương pháp cổ điển phụ thuộc mạnh vào (i) chất lượng tiền xử lý, (ii) độ đúng của mQRS/fQRS detection, và (iii) giả định mô hình (tuyến tính/phi tuyến, độc lập nguồn, phân bố Gauss). Trong bối cảnh fECG có SNR thấp và hay chồng lấp với mECG, các giả định này dễ bị vi phạm, dẫn tới kết quả kém vững và khó tổng quát khi triển khai ngoài môi trường kiểm soát.


\section{Nhóm phương pháp học sâu}
\subsection{Nhóm mô hình CNN cho fQRS detection và/hoặc trích dạng sóng fECG}

\begin{figure}[h]
  \centering
  \includegraphics[width=0.8\textwidth]{Hinhve/image_41.png}
  \caption{Sơ đồ khối kiến trúc của mạng CNN.}
  \label{fig:example_image}

\end{figure}

CNN được sử dụng rộng rãi nhờ khả năng bắt các mẫu cục bộ (đỉnh, sườn dốc, cấu trúc QRS) và xử lý chuỗi thời gian thô hoặc đã tiền xử lý tương đối hiệu quả. Các hướng triển khai tiêu biểu gồm: (i) dùng STFT để tạo biểu diễn thời--tần rồi đưa vào CNN phục vụ phát hiện fQRS; (ii) CNN 1D sâu hơn kết hợp phân đoạn cửa sổ và gán nhãn vùng chứa fQRS; (iii) tăng bền vững với nhiễu bằng cách đưa đặc trưng đánh giá nhiễu như SampEn vào pipeline; (iv) các biến thể ResNet/Octave Convolution nhằm giảm dư thừa và tăng hiệu quả tính toán; (v) các thiết kế kết hợp miền phổ--wavelet và CNN, hoặc cơ chế nhất quán miền (cross-domain consistency) để cải thiện khả năng tổng quát hoá. Nhìn chung, CNN có ưu thế về hiệu quả và tiềm năng thời gian thực, nhưng hạn chế thường gặp là “cửa sổ tiếp nhận” (receptive field) hữu hạn và phụ thuộc cấu hình, khiến việc nắm quan hệ dài hạn hoặc biến thiên nhịp--hình thái theo thời gian có thể chưa tối ưu nếu chỉ dùng CNN thuần
\cite{lee_fetal_2018, cao_morphology_2022, zhong_deep_2018, vo_efficient_2020}.


\subsection{Nhóm mô hình RNN (LSTM/GRU) cho phụ thuộc thời gian} 
RNN được thiết kế cho dữ liệu chuỗi nên phù hợp để mô hình hoá phụ thuộc theo thời gian và biến thiên nhịp giữa các chu kỳ tim. LSTM/GRU giúp khắc phục vấn đề tiêu biến gradient của RNN truyền thống và có thể khai thác ngữ cảnh dài hơn để ổn định phát hiện nhịp hoặc tinh chỉnh đầu ra theo chuỗi. Khi áp dụng trực tiếp lên tín hiệu bụng thô, trạng thái ẩn có thể bị chi phối bởi mQRS do biên độ mECG lớn, làm giảm nhạy với các spike fQRS yếu trong điều kiện SNR thấp; vì vậy nhiều công trình dùng RNN như một tầng tinh chỉnh sau khi đã có đặc trưng cục bộ tốt
\cite{hochreiter_long_1997, darmawahyuni_accurate_2023, darmawahyuni_deep_2023}.


\subsection{Nhóm mô hình lai CNN--RNN}
Hybrid CNN--RNN tận dụng ưu điểm bổ sung: CNN trích đặc trưng cục bộ liên quan đến hình thái (đặc biệt QRS), còn RNN học phụ thuộc thời gian và tính không dừng của tín hiệu. Một cấu trúc điển hình là CNN ở các tầng đầu để làm sạch/nhấn mạnh đặc trưng fQRS, sau đó GRU/LSTM/BiLSTM để khai thác sự liên tục theo thời gian và giảm nhiễu quyết định. Một số công trình còn kết hợp DWT nhằm khử nhiễu trước khi vào mạng. Nhóm hybrid thường đạt hiệu năng tốt hơn trong môi trường nhiễu hoặc dữ liệu khác miền, nhưng đổi lại chi phí tính toán và độ phức tạp huấn luyện tăng \cite{ghonchi_common_2022,darmawahyuni_accurate_2023, darmawahyuni_deep_2023}.

\subsection{Nhóm mô hình sinh (GAN/CycleGAN) hướng tới khôi phục dạng sóng và bảo toàn morphology}
Các kiến trúc sinh như GAN/CycleGAN được quan tâm vì mục tiêu không chỉ “bắt đỉnh fQRS” mà còn khôi phục toàn bộ dạng sóng fECG với hình thái P--QRS--T (phục vụ phân tích khoảng PR/ST/QT và các bất thường hình thái). CycleGAN cho phép huấn luyện trong bối cảnh không ghép cặp (unpaired) bằng cơ chế cycle-consistency, giúp mô hình học ánh xạ giữa miền tín hiệu bụng và miền fECG đồng thời hạn chế mất chi tiết hình thái. Các biến thể đưa attention hoặc ràng buộc tương quan nhằm tăng khả năng “nhìn thấy” các thành phần fECG bị che bởi mECG và nhiễu. Nhìn chung, nhóm sinh có lợi thế về chất lượng tái tạo và bảo toàn chi tiết, nhưng huấn luyện khó hơn, nhạy hyperparameter và thường nặng hơn về tài nguyên. 
\cite{goodfellow_generative_2014, zhu_unpaired_2017, mohebbian_fetal_2022, shokouhmand_fetal_2023}.

\subsection{Nhóm Autoencoder và họ U-Net/UNet++/UNet3+ cho bài toán trích dạng sóng fECG}
Autoencoder (AE) và các biến thể encoder--decoder là lựa chọn tự nhiên cho bài toán “tách/khôi phục” tín hiệu: encoder nén và học biểu diễn, decoder tái tạo đầu ra; skip connection giúp giữ thông tin cục bộ và ổn định gradient, từ đó cải thiện bảo toàn hình thái theo nhịp. Survey cho thấy nhiều hướng phát triển dựa trên U-Net và các biến thể (UNet++, UNet3+) để tăng hợp nhất đa tỉ lệ và thu hẹp “semantic gap” giữa encoder--decoder. Ngoài ra, có các hướng tách thành phần mẹ--thai bằng kiến trúc kép (ví dụ dual Res-UNet) hoặc tích hợp nhiều module để tiến tới một lần suy luận (single-pass) thay vì pipeline nhiều bước. Một nhánh thực dụng khác là LinkNet/LinkNet++ với residual/dense blocks và cơ chế merge đặc trưng hiệu quả, nhằm cân bằng giữa chất lượng tái tạo và chi phí tính toán.

\section{Nhận xét, khoảng trống nghiên cứu và định vị đề tài}
% TODO: viết thành văn xuôi (gạch đầu dòng định hướng) — phần "ăn điểm" định vị đóng góp
\begin{itemize}
  \item Tổng kết ưu/nhược: cổ điển phụ thuộc tiền xử lý + giả định mô hình; học sâu mạnh hơn nhưng nặng/khó huấn luyện.
  \item \textbf{W-NETR là SOTA hiện tại} cho tách fECG dạng sóng (hai nhánh UNETR + khử đặc trưng mECG).
  \item \textbf{Khoảng trống:} khối MLP-FFN trong Transformer dùng kích hoạt phi tuyến cố định $\Rightarrow$ hạn chế biểu diễn/bảo toàn hình thái.
  \item \textbf{Định vị đề tài:} thay MLP-FFN bằng KAN (kích hoạt học được theo spline) $\Rightarrow$ động lực cho mô hình đề xuất ở Chương~\ref{chapter:ProposedModel}.
\end{itemize}

\section{Kết luận chương}
% TODO: viết thành văn xuôi (gạch đầu dòng định hướng)
\begin{itemize}
  \item Tóm tắt bức tranh phương pháp và lý do chọn hướng KAN; chuyển ý sang cơ sở lý thuyết học sâu (Chương~\ref{chapter:DLBackground}).
\end{itemize}

\end{document}
